\section{Introduction}\label{sec:intro}

\subsection{The Kessler Challenge and Solar Cycle 25}
The modern LEO environment is defined by two concurrent trends: sustained mega-constellation expansion and heightened geomagnetic volatility near the Solar Cycle 25 climax (2025--2026). The resulting traffic density amplifies collision cascades first formalized in the Kessler framework \cite{kessler1978collision}, while recent constellation studies show that propagation error now scales into direct operational risk for daily conjunction screening \cite{chen2025oneweb,kang2025solar}. During storm-time conditions, thermospheric heating drives rapid drag increases and measurable orbital decay across LEO shells \cite{oliveira2021current,walter2025orbital}.

\subsection{Limits of Traditional and Pure ML Propagators}
SGP4 remains the operational default because of speed and interoperability, but its pseudo-ballistic $B^*$ representation is effectively piecewise static between TLE updates. Under severe geomagnetic forcing, that assumption breaks down and miss-distance growth accelerates, as reported during the May 2024 G5 event \cite{parker2024satellite,dey2024improving}. Purely data-driven residual correctors improve short-horizon performance in quiet regimes but frequently violate initial-state consistency: a network-predicted residual at $t=0$ causes non-physical displacement at the TLE epoch. We refer to this failure mode as \textbf{Phantom Drift}, consistent with instability patterns observed in unconstrained orbital ML pipelines \cite{kozhaya2021comparison,bertolini2024hybrid}.

\subsection{The Gated PINN Solution}
To address both storm resilience and kinematic consistency, we introduce a Gated Physics-Informed Neural Network (G-PINN). The architecture applies a differentiable Physics Gate,
\begin{equation}
\Gamma(t)=\tanh(\lambda t),
\end{equation}
as a hard boundary condition on residual outputs, guaranteeing zero correction at $t=0$ and therefore zero epoch drift. The model fuses kinematic state variables with space-weather drivers ($K_p$, $F_{10.7}$), then deploys corrected trajectories inside a probabilistic conjunction stack. This framework aligns with current PINN advances in orbit determination and uncertainty-aware space safety systems \cite{loshelder2025od,ramesh2025space}.

\subsection{Contributions and Paper Structure}
This paper contributes: (i) a hard-constrained gate that eliminates Phantom Drift, (ii) storm-aware residual learning under Solar Cycle 25 conditions, (iii) model-driven discovery of a persistent $+5.5$ m/s radial bias in baseline SGP4 handling, and (iv) an operational Monte Carlo risk pipeline with Mahalanobis classification and WebGL visualization. The remainder of the paper is organized as follows: Section~\ref{sec:litreview} reviews prior work, Section~\ref{sec:method} details the G-PINN framework, Sections~\ref{sec:results} and \ref{sec:ops} present scientific and operational findings, and Section~\ref{sec:conclusion} concludes with forward directions.